\documentclass[11pt]{article}
\usepackage[T1]{fontenc}
\usepackage{textcomp}
\usepackage[margin=0.75in]{geometry}
\usepackage{amsmath,amssymb,amsthm}
\usepackage{array,booktabs}
\usepackage{hyperref}
\setlength{\parindent}{0pt}
\newtheorem{theorem}{Theorem}
\theoremstyle{definition}
\newtheorem{definition}{Definition}
\title{Flow Proof Artifact: List-len-append-undec-singleton}
\author{Generated by \texttt{flow doc proof}}
\date{Flow Proof Book}
\begin{document}
\maketitle
Appending eleven singletons gives length eleven.\\[0.5em]
\textbf{Source.} church — https://en.wikipedia.org/wiki/Length\_of\_a\_list\\[0.5em]
\bigskip
\noindent{\large \textbf{Derived fact 1.} \textit{len of eleven appended singletons equals 11} \hfill List \cdot length \cdot append of eleven singletons has length eleven}\par
\smallskip\noindent\textit{Coordinate: List \cdot length \cdot append of eleven singletons has length eleven}\par
\smallskip\noindent\textit{Source: church}\par
\smallskip\noindent\textit{Built on: append preserves length sum, for length on List, append of ten singletons has length ten, for length on List}\par
\medskip
\noindent\textbf{Goal.} len of eleven appended singletons equals 11\par
\noindent$\displaystyle \forall a \in \mathbb{N} \forall b \in \mathbb{N} \forall c \in \mathbb{N} \forall d \in \mathbb{N} \forall e \in \mathbb{N} \forall f \in \mathbb{N} \forall g \in \mathbb{N} \forall h \in \mathbb{N} \forall i \in \mathbb{N} \forall j \in \mathbb{N} \forall k \in \mathbb{N}\quad len(cons(a, nil) ++ cons(b, nil) ++ cons(c, nil) ++ cons(d, nil) ++ cons(e, nil) ++ cons(f, nil) ++ cons(g, nil) ++ cons(h, nil) ++ cons(i, nil) ++ cons(j, nil) ++ cons(k, nil)) = 11$\par
\medskip
\noindent
\begin{tabular}{@{} >{\raggedright\arraybackslash}p{0.06\textwidth} >{\raggedright\arraybackslash}p{0.47\textwidth} >{\raggedleft\arraybackslash}p{0.41\textwidth} @{}}
\textbf{\#} & \textbf{Proof} & \textbf{Mathematics} \\
\midrule
\textbf{1.} & We prove that append of eleven singletons has length eleven for length on List. & \\[0.45em]
\textbf{2.} & We invoke the derived fact governing length on List: append preserves length sum, for length on List (instantiated for cons(a, nil) ++ cons(b, nil) ++ cons(c, nil) ++ cons(d, nil) ++ cons(e, nil) ++ cons(f, nil) ++ cons(g, nil) ++ cons(h, nil) ++ cons(i, nil) ++ cons(j, nil), cons(k, nil)). & \\[0.45em]
\textbf{3.} & We invoke the derived fact governing length on List: append of ten singletons has length ten, for length on List (instantiated for a, b, c, d, e, f, g, h, i, j). & \\[0.45em]
\textbf{4.} & From step 2 and step 3, this implies len(cons(a, nil) plus plus cons(b, nil) plus plus cons(c, nil) plus plus cons(d, nil) plus plus cons(e, nil) plus plus cons(f, nil) plus plus cons(g, nil) plus plus cons(h, nil) plus plus cons(i, nil) plus plus cons(j, nil) plus plus cons(k, nil)) equals 11. Hence proven. & $\displaystyle len(cons(a, nil) ++ cons(b, nil) ++ cons(c, nil) ++ cons(d, nil) ++ cons(e, nil) ++ cons(f, nil) ++ cons(g, nil) ++ cons(h, nil) ++ cons(i, nil) ++ cons(j, nil) ++ cons(k, nil)) = 11$ \\[0.45em]
\bottomrule
\end{tabular}
\label{thm:List-length-append_of_eleven_singletons_has_length_eleven}
\par\medskip\hrule
\end{document}
